\documentclass[11pt,a4paper,sans]{moderncv}
\moderncvstyle{banking}
\moderncvcolor{black}
\nopagenumbers{}
\usepackage[utf8]{inputenc}
\usepackage{ragged2e}
\usepackage[scale=0.915]{geometry}
\usepackage{import}
\usepackage{multicol}
\usepackage{enumitem}
\usepackage{amssymb}
\name{Elio}{Govea}
\newcommand*{\customcventry}[7][.13em]{
\begin{tabular}{@{}l}
{\bfseries #4} \
{\itshape #3}
\end{tabular}
\hfill
\begin{tabular}{l@{}}
{\bfseries #5} \
{\itshape #2}
\end{tabular}
\ifx&#7&%
\else{\
\begin{minipage}{\maincolumnwidth}%
\small#7%
\end{minipage}}\fi%
\par\addvspace{#1}}
\begin{document}
\makecvtitle
\vspace*{-16mm}
\begin{center}\textbf{Software Developer}\end{center}
\begin{center}
\begin{tabular}{ c c c }
\faMobile \enspace eliogovea1993@gmail.com \enspace & \faHome\enspace Helsinki, Finland \\
\faLinkedin\enspace \color{blue} \href{https://www.linkedin.com/in/elio-alejandro-govea-aguilar-557598172/}{elio-alejandro-govea-aguilar} \enspace &
\faGithub\enspace \color{blue} \href{https://github.com/eliogovea}{eliogovea}
\end{tabular}
\end{center}

\section{Profile}
Software developer. Currently building aviation weather systems at {\color{blue}\href{https://www.vaisala.com/en}{Vaisala}}. Former {\color{blue}\href{https://icpc.global/ICPCID/PNF1J4PZOUVQ}{ICPC}} World Finalist

\section{Experience}

\customcventry
{02/2025 ‐ present}
{{\color{blue}\href{https://www.vaisala.com/en}{Vaisala}}}
{Senior Software Engineer}
{Vantaa, Finland}
{}
{}

\begin{itemize}[leftmargin=0.6cm, label={\textbullet}]
    \item Contributing to {\color{blue}\href{https://www.vaisala.com/en/products/systems/avimet}{AviMet}}, an aviation weather information system used at airports worldwide --- providing real-time wind, visibility, lightning, wind shear, and runway-surface data to air-traffic control and pilots
    \item Working across a wide variety of issues, implementing new features, and adding support for new platforms
\end{itemize}

\customcventry
{05/2023 ‐ 02/2025}
{{\color{blue}\href{https://www.f-secure.com/}{F-Secure}}}
{Senior Developer}
{Helsinki, Finland}
{}
{}

\begin{itemize}[leftmargin=0.6cm, label={\textbullet}]
    \item Designed CDN cost-reduction strategy for the Sense SDK --- cut delivery costs by \emph{80\%} with minimal impact on end-users
        \begin{itemize}
            \item Analyzed CDN usage and identified key optimization areas
            \item Designed a custom caching strategy to significantly decrease the volume of data transferred via CDN
            \item Ensured that changes were backend-focused to avoid client-side updates and lengthy deployment negotiations
        \end{itemize}
\end{itemize}

\customcventry
{12/2021 ‐ 04/2023}
{{\color{blue}\href{https://www.f-secure.com/}{F-Secure}}}
{Developer}
{Helsinki, Finland}
{}
{}

\begin{itemize}[leftmargin=0.6cm, label={\textbullet}]
    \item \textbf{Built Sense SDK} --- F-Secure's C++ SDK for home-router security: URL filtering and parental controls (forcing safe-mode browser engines), traffic-statistics collection, device recognition and fingerprinting; across IPv4/IPv6, DHCP, TCP/UDP, DNS, HTTP/HTTPS using Linux networking interfaces (netlink, netfilter, iptables, nfqueue, nflog, conntrack, BPF) and integrated with cloud back-end services
    \item \textbf{Shipped on x86, ARM, and MIPS targets} --- designed and optimized the cross-platform build system (CMake + GNU Make) for fast, reproducible builds across toolchains, plus SDK integrations for reference platforms to streamline demos and future integrations
    \item \textbf{Designed Jenkins / AWS CI pipelines} --- enforced code quality through coverage, static, and dynamic analysis; produced builds for multiple architectures and reference platforms
    \item \textbf{Built testing infrastructure in Python and Robot Framework}; implemented recovery mechanisms for OS- and cloud-service-level errors
    % \item Notes
    %     \begin{itemize}[leftmargin=0.6cm, label={\textbullet}]
    %         \item C | C++ | Python | Project structure
    %         \item ARM | MIPS
    %         \item Linux
    %             \begin{itemize}
    %                 \item sh | socket | netlink | netfilter | conntrack | nfqueue | nflog | iptables
    %                 \item OpenWrt | reference platforms | easy deployment | simpler demo | packages (.ipk) | firmware 
    %                 \item buildroot | GCC | cross compile toolchains
    %             \end{itemize}
    %         \item Netowrking
    %             \begin{itemize}
    %                 \item IPv4 | IPv6 | UDP | TCP | DNS | HTTP | HTTPS
    %             \end{itemize}
    %         \item clients for cloud services (registration) | localization
    %     \end{itemize}
\end{itemize}

\customcventry{09/2020 ‐ 11/2021}
{{\color{blue}\href{https://www.f-secure.com/}{F-Secure}}}
{Software Developer}
{Poznań, Poland}
{}
{}

\begin{itemize}[leftmargin=0.6cm, label={\textbullet}]
    \item Contributed to the core C++ of the Sense SDK
    \item Optimized a notifications system for security events
    \item Developed embedded Linux clients to consume cloud services
    \item Added endpoints to existing REST APIs
    \item Improved internal C++ interfaces and data structures
    \item Fixed bugs on existing products
    \item Added new interfaces using Unix domain socket for communication with different client/services
    % \item Notes
    %     \begin{itemize}
    %         \item signals | reduce parsing | unit tests | traffic statistics | UDP traffic filtering | TCP traffic filtering | threads sychronization | utilities to monitor services running on different threads | implemented interfaces using unix domain sockets to communicate with different client/services | implemented C++ utilities to handle multiple OS operations | add testing features
    %     \end{itemize}
\end{itemize}


\customcventry{01/2018 ‐ 08/2020}{\color{blue}\href{https://www.upr.edu.cu/}
{Universidad de Pinar del Rio}}
{Software Engineer}
{Pinar del Rio, Cuba}
{}
{}

\begin{itemize}[leftmargin=0.6cm, label={\textbullet}]
    \item Designed and implemented applications to collect and visualize vibration signals using C++ and Qt
        \begin{itemize}
            \item Real-time signal processing
            \item Technologies used include: C++ | IMU | Arduino | Serial communication | Qt
        \end{itemize}
    \item Designed and implemented applications to collect and visualize data from sensors for object position tracking
        \begin{itemize}
            \item Technologies used include: C++ | IMU | Arduino | Serial communication | Linux | OpenGL
        \end{itemize}
    \item Assisted in teaching programming courses
        \begin{itemize}
            \item Prepared lectures on C++, Python, Arduino, Linux, Raspberry Pi
            \item Conducted lab sessions and provided assistance to students with programming projects 
            \item Designed and graded assignments
        \end{itemize}
    \item ICPC Judge for the Caribbean region
        \begin{itemize}
            \item Evaluated and proposed challenging programming problems and solution for the ICPC Caribbean Finals
            % \item Collaborated with fellow judges to maintain high competition standards and ensure a fair and engaging experience
        \end{itemize}
    \item ICPC coach
        \begin{itemize}
            \item Coached a team to the ICPC Latin American Finals
        \end{itemize}
\end{itemize}

\section{Education}
\customcventry{09/2012-07/2017}{\color{blue}\href{https://www.upr.edu.cu/}{UPR}}
{B.Sc Telecommunications and Electronics Engineering}
{Pinar del Rio, Cuba}
{}
{}
{}

\begin{itemize}[leftmargin=0.6cm, label={\textbullet}]
    \item GPA \emph{4.95/5.00}
    \item ICPC world finalist
\end{itemize}

\section{Achievements}

{\begin{itemize}[label=\textbullet]
\item {\textbf{Honourable Mention} 2018 ACM ICPC World Finals}
\item {\textbf{2nd Place (Silver Medal)} 2016 ACM ICPC Latin 
American Finals}
\item {\textbf{4th Place} 2015 ACM ICPC Latin 
American Finals}
\item {\textbf{Top 1000} 2019 Google Code Jam}
\end{itemize}}

\section{Projects}

\begin{itemize}[label=\textbullet]
    \item \href{https://github.com/eliogovea/baseball-limits-2d}{\textbf{Baseball Limits 2D}} (\href{https://eliogovea.github.io/baseball-limits-2d/}{\color{blue}live demo})
        \begin{itemize}
            \item What does ``best ever'' look like when you compare two baseball stats at once? In one dimension it's a ranking; in two it's a Pareto frontier --- the players no one strictly outperforms in both
            \item Interactive viewer over 150 years of MLB data (Lahman Database 1871--2023, Baseball Reference 2024)
        \end{itemize}
    \item \href{https://github.com/eliogovea/dev-env-cpp}{\textbf{C++ project template}}
        \begin{itemize}
            \item Docker environment with common tools
            \item CMake presets for multiple configurations
            \item CI pipelines for reporting code coverage and generating documentation
        \end{itemize}
    \item \href{https://github.com/eliogovea/dns}{\textbf{DNS library}}
        \begin{itemize}
            \item Parse and generate DNS messages
            \item Provide examples of DNS resolution
        \end{itemize}
    \item \href{https://github.com/eliogovea/suffix-automaton-vis}{\textbf{Suffix automaton visualization}} (\href{https://eliogovea.github.io/suffix-automaton-vis/}{\color{blue}live demo})
        \begin{itemize}
            \item Interactive visualization tool using TypeScript and D3 to illustrate the algorithm to build a suffix automaton
        \end{itemize}
    \item \href{https://github.com/eliogovea/compile-time-cpp}{\textbf{Template metaprogramming examples}}
        \begin{itemize}
            \item C++ examples to create and modify types and constants at compile time
        \end{itemize}
    \item \href{https://github.com/eliogovea/router-network-docker}{\textbf{Router-network Docker testbed}}
        \begin{itemize}
            \item Docker setup that simulates a small home network --- clients behind a router-as-gateway talking to remote servers --- for testing network applications end-to-end without real hardware
        \end{itemize}
    \item \textbf{Competitive programming archive}: \href{https://github.com/eliogovea/icpc-reference}{\color{blue}reference} \& \href{https://github.com/eliogovea/solutions_cp}{\color{blue}solutions}
        \begin{itemize}
            \item Personal notebook of algorithms and data structures, the kind ICPC teams build up over years
            \item Running set of solutions to online-judge problems
        \end{itemize}
\end{itemize}

\section{Skills}

{\begin{itemize}[label=\textbullet]
\item \textbf{Languages:} C | C++ | Python | Bash
\item \textbf{Build \& tooling:} CMake | GNU Make | Git | Docker | LXC
\item \textbf{Systems:} Linux | Buildroot | OpenWrt
\item \textbf{Concepts:} Algorithms | Data Structures | Networking
\item \textbf{Protocols:} IP | DHCP | UDP | TCP | DNS | HTTP | TLS | QUIC
\end{itemize}}

\section{Languages}

\begin{multicols}{2}
    \begin{itemize}[label=\textbullet]
    \item \textbf{Spanish} [Native]
    \item {\textbf{English} [Full Professional Proficiency]}
    \end{itemize}
\end{multicols}
\end{document}